\begin{abstract}
\noindent Synthetic respondents built from large language models promise a cheap first reading of survey items before fielding, yet their value as a screening instrument is rarely tested on items they were not built for. We evaluate synthetic cultural agents, which compile the signs of declared country-level preference benchmarks into parameter adapters through a regularized pairwise choice model, as a scoring instrument for new item text. Under a preregistered protocol, the adapters alone assigned a dimension and a direction to candidate World Values Survey items before any human value was read. Against the existing survey, the selected items predict the direction of cross-national differences better than random assignment, but the signal is weak and no preference dimension passes both a second-bank check and the held-out screen. A factorial over sign profiles shows that the dimensions are entangled, more so under prompting. Because the human criterion is retrospective, pre-piloting remains a prospective use that this paper does not test.
\end{abstract}

\noindent\textbf{Keywords:} synthetic respondents; large language models; survey item screening; preregistration; cross-national measurement; direct preference optimization.
